\chapter{局部随机神经网络离散方法}
\section{随机特征与物理导数}
每侧用$N_i$个固定特征$\phi_{ij}$，令$u_i^N=\sum_{j=1}^{N_i}c_{ij}\phi_{ij}$。隐藏权重与偏置随机生成后不优化。直线子域先映射到单位参考方形，左/右侧$\xi=(x-x_{\min})/.5,\eta=y$，特征为
\begin{equation}
\phi_j=\tanh\{w_{j1}(\xi-.5)+w_{j2}(\eta-.5)+b_j\}.
\end{equation}
若$t=\tanh z$，则$\sigma'=1-t^2$、$\sigma''=-2t(1-t^2)$。物理梯度带$x$方向因子2，物理Laplacian为$\sigma''(4w_{j1}^2+w_{j2}^2)$。忽略局部坐标尺度会产生模型错误，故程序以独立中心差分检查链式导数。

隐藏参数取指定均匀范围，不假设不同宽度的特征严格嵌套。相同种子改变$N_i$时，权重和偏置的随机调用顺序可能不同；宽度实验因此是同分布统计比较，不构成嵌套空间的单调性证明。

\section{分块矩阵组装}
令$c=(c_1^T,c_2^T)^T$。内部行分别为$(-\beta_1\Delta\Phi_1,0)$与$(0,-\beta_2\Delta\Phi_2)$，外边界仅激活所属侧的特征块。界面两类行是
\begin{equation}
A_u=(-\Phi_1,\Phi_2),\qquad
A_q=(-\beta_1\partial_n\Phi_1,\beta_2\partial_n\Phi_2).
\end{equation}
通量必须含各侧$\beta$，不能用导数连续代替通量条件。界面端点或角点可能造成多组约束重合，初期采用段中点或开区间随机点以避免重复。

各组样本数为$m_k$、平方权重为$\lambda_k$、固定物理尺度为$s_k$时，目标写为
\begin{equation}\label{eq:weighted}
\min_c\sum_k\frac{\lambda_k}{m_ks_k^2}\norm{A_kc-d_k}_2^2.
\end{equation}
叠置时矩阵和数据同时乘$\sqrt{\lambda_k/m_k}/s_k$。前期直线模型取边界与界面平方权重50；raw尺度全1，balanced内部尺度$\beta_i$、通量尺度$\beta_{\max}$。这些选择改变经验目标中的相对重要性，不是无损等价变换。高对比度时归一化通量小仍可对应较大物理通量误差，所以两种指标都需报告。

\subsection{分组目标与权重的含义}
设内部残差向量为$r_1,r_2$，每侧边界向量为$r_{B,i}$，两类界面向量为$r_u,r_q$。直线问题的目标可具体写为
\begin{align}
J(c)={}&\frac{\norm{r_1}^2}{m_1s_1^2}+\frac{\norm{r_2}^2}{m_2s_2^2}
+\sum_{i=1}^2\frac{\lambda_B\norm{r_{B,i}}^2}{m_{B,i}}\nonumber\\
&+\frac{\lambda_\Gamma\norm{r_u}^2}{m_\Gamma}
+\frac{\lambda_\Gamma\norm{r_q}^2}{m_\Gamma s_q^2}.
\end{align}
代码把每侧边界作为同点数的独立组归一化，两组的合计贡献在不同策略间保持一致。$m_\Gamma$同时进入两类界面行，不能把界面点数增加误解为界面平方权重也按同一比例增加。

行乘数$\gamma$进入平方目标后是$\gamma^2$，因此花形近似复现的行乘数50与直线模型的平方权重50不是同一个设置。不同残差可能对应不同物理单位；本文几何坐标无量纲，尺度选择是明确的数值约定，尚未建立跨物理模型的普适量纲平衡原则。正文同时报告物理通量，是为了防止归一化数值遮蔽实际界面偏差。

\subsection{约束数量与近似空间}
通常需要足够行数支撑$K$个系数，但$m_{\rm eq}\ge K$只是满列秩的必要条件，不是充分条件。随机特征近相关、某侧覆盖不足或重复约束都可能使矩阵近秩亏。另一方面，增加点数没有改变固定特征空间本身，不能创造网络缺少的局部尺度。新增约束如何改变解及条件数还取决于误差分布、权重和截断，见命题\ref{prop:rows}。

\section{求解与独立评估}
采用SVD型最小二乘避免显式正规方程，报告未截断条件数、保留条件数、有效秩、最小保留奇异值及系数范数。非零岭参数通过增广$[B;\sqrt\alpha I]$求解，数据矩阵秩与增广求解秩分开记录。

程序正常且指标有限仅说明运行完成，不代表精度验收或理论收敛。解误差在独立点上计算；均匀等权点的相对离散误差为$\sqrt{\sum|u_N-u|^2/\sum|u|^2}$，并报告最大误差。PDE残差使用解析算子在独立点计算。训练损失可能很小而边界、界面或未采区域误差仍大，有限点零残差不能保证连续PDE被满足。

命题\ref{prop:finite_error}给出一个解释框架：小训练目标只反映已选行及其权重，并不自动控制验证近似项、截断丢弃项或预测映射。本文测试点与训练点分离、另报物理通量及秩，正是为了降低只按一个训练数字判断算法的风险。验证点也只构成离散观测，不能代替完整连续误差分析。

\section{有限差分导数模式与原文区别}
对每个子域的平滑特征采用
\begin{align}
\partial_x\phi&\approx\frac{\phi(x+h_1,y)-\phi(x-h_1,y)}{2h_1},\\
\partial_{xx}\phi&\approx\frac{\phi(x+h_2,y)-2\phi(x,y)+\phi(x-h_2,y)}{h_2^2}.
\end{align}
即使点靠近界面，也只对同一网络特征延拓求差分，不切换到另一侧网络。对足够光滑函数，两种中心格式的截断误差为$O(h^2)$；二阶差分中浮点相减误差会除以$h_2^2$放大，因此步长并非越小越好。

花形近似复现按原文使用物理坐标$z=w\cdot x+b$及行权重$\gamma=50$\cite{interfacepreprint}，没有沿用直线模型的参考坐标、样本数归一化或平方权重50。其初始化及积分细节见第五章。本文强形式实现不是DG或PG方法，也不将解析导数版本称为原文FD严格复现。
